\documentclass[a4paper,12pt]{article}
\usepackage[utf8]{inputenc}
\usepackage[T1]{fontenc}
\usepackage[ngerman]{babel}
\usepackage{amsmath}
\usepackage{amssymb}
\usepackage{enumitem}
\usepackage{geometry}
\geometry{a4paper, margin=2.5cm}

\title{Einsendeaufgaben -- Aufgabe 4.4\\[0.5em]
\large 61211 Analysis}
\author{}
\date{}

\begin{document}
\maketitle

\section*{Aufgabe 1}

\begin{enumerate}[label=\alph*)]

    \item
    Sei $r > 0$ beliebig.
    \begin{align*}
        D_1 f(x,y,z) &= \frac{1}{x} \\
        D_2 f(x,y,z) &= \frac{2}{y} \\
        D_3 f(x,y,z) &= \frac{3}{z}
    \end{align*}
    \[
        g(x,y,z) = x^2 + y^2 + z^2 - 6r^2 \qquad M(g) := \{z \in M \mid g(z) = 0\}
    \]
    \begin{align*}
        D_1 g(x,y,z) &= 2x \\
        D_2 g(x,y,z) &= 2y \\
        D_3 g(x,y,z) &= 2z
    \end{align*}
    \[
        g'(x,y,z) = \begin{pmatrix} 2 & 0 & 0 \\ 0 & 2 & 0 \\ 0 & 0 & 2 \end{pmatrix}
    \]
    Für jedes $z \in M(g)$ hat $g'(z)$ also den Höchstrang. Wir können demnach den Satz 4.5.9 über Extrema unter Nebenbedingungen anwenden.

    $f'(c) = \lambda g'(c)$ ist also eine notwendige Bedingung dafür, dass $c \in M(g)$ ein Extremum ist.

    Aufstellung der Gleichungen:
    \[
        \frac{1}{x} = \lambda 2x \qquad \frac{2}{y} = \lambda 2y \qquad \frac{3}{z} = \lambda 2z
    \]
    \[
        x^2 + y^2 + z^2 = 6r^2
    \]
    Umstellung nach $x^2$, $y^2$ und $z^2$:
    \begin{align*}
        x^2 &= \frac{1}{2\lambda} \\
        y^2 &= \frac{1}{\lambda} \\
        z^2 &= \frac{3}{2\lambda}
    \end{align*}
    Einsetzen von $x^2$, $y^2$ und $z^2$ in $x^2 + y^2 + z^2 = 6r^2$:
    \[
        \frac{1}{2\lambda} + \frac{1}{\lambda} + \frac{3}{2\lambda} = 6r^2
    \]
    \[
        \Rightarrow \frac{3}{\lambda} = 6r^2
    \]
    \[
        \Rightarrow \lambda = \frac{1}{2r^2}
    \]
    Nun Bestimmung von $x$:
    \[
        x^2 = \frac{1}{2\lambda}
    \]
    \[
        \Rightarrow x^2 = r^2
    \]
    Da $x > 0$ und $r > 0$ folgt
    \[
        x = r
    \]
    Für $y, z$ erhalten wir genauso
    \begin{align*}
        y &= \sqrt{2}\, r \\
        z &= \sqrt{3}\, r
    \end{align*}
    Das heißt jedes lokales Extremum $c \in M(g)$ von $f|_{M(g)}$ hat die Form $c = (r, \sqrt{2}\, r, \sqrt{3}\, r)$.

    Wir zeigen, dass $c := (r, \sqrt{2}\, r, \sqrt{3}\, r)$ ein lokales Maximum ist.
    \begin{align*}
        D_{11} f(x,y,z) &= -\frac{1}{x^2} & D_{12} f(x,y,z) &= D_{21} f(x,y,z) = 0 \\
        D_{22} f(x,y,z) &= -\frac{2}{y^2} & D_{13} f(x,y,z) &= D_{31} f(x,y,z) = 0 \\
        D_{33} f(x,y,z) &= -\frac{3}{z^2} & D_{23} f(x,y,z) &= D_{32} f(x,y,z) = 0
    \end{align*}
    \[
        H(r, \sqrt{2}\, r, \sqrt{3}\, r) = -\frac{1}{r^2} I_3
    \]
    \[
        Q(u) = -\frac{1}{r^2}(u_1^2 + u_2^2 + u_3^2) < 0 \quad \text{für jedes } u \in \mathbb{R}.
    \]
    Demnach ist $c = (r, \sqrt{2}\, r, \sqrt{3}\, r)$ ein lokales Maximum unter der Nebenbedingung $g$. $c$ ist also auch nach Satz 4.5.9 das einzige lokale Extremum von $f|_{M(g)}$.

    Noch zu zeigen, dass $c$ das globale Maximum von $f|_{M(g)}$ ist. Jeder Grenzwert $a$ einer konvergenten Folge $(a_n)$ mit $g(a_n) = 0$ für alle $n \in \mathbb{N}$ hat die Eigenschaft $a \in M(g)$, da aus der Stetigkeit von $g$
    \[
        0 = \lim_{a_n \to a} g(a_n) = g(a)
    \]
    folgt. Also ist $M(g)$ abgeschlossen. $M(g)$ ist ebenfalls beschränkt nach der euklidischen Norm, da die Nebenbedingung äquivalent zu $\sqrt{x^2 + y^2 + z^2} = \sqrt{6}\, r$ ist. $M(g)$ ist nach der euklidischen Norm die Oberfläche der Kugel mit dem Radius $\sqrt{6}\, r$. Aus der Abgeschlossenheit und Beschränkung folgt, dass $M(g)$ kompakt ist. Wegen der Kompaktheit von $M(g)$ und der Stetigkeit von f ist $c$ also das Maximum von $f|_{M(g)}$.

    \item
    Wir können $a, b$ und $c$ mit $a, b, c > 0$ wie folgt definieren:
    \[
        a := x^2, \quad b := y^2 \text{ und } c := z^2 \text{ mit } x, y, z > 0.
    \]
    Außerdem
    \[
        r := \sqrt{\frac{x^2 + y^2 + z^2}{6}}
    \]
    Dann folgt
    \begin{align*}
        \ln(a b^2 c^3) &= 2 \ln(x y^2 z^3) \\
        &= 2(\ln(x) + 2 \ln(y) + 3 \ln(z)) \\
        &= 2 f(x,y,z) \\
        &\leq 2 f(r, \sqrt{2}\, r, \sqrt{3}\, r) \quad \text{(wegen } 6r^2 = x^2 + y^2 + z^2 \text{ können wir a) anwenden)} \\
        &= 2(\ln(r) + 2 \ln(\sqrt{2}\, r) + 3 \ln(\sqrt{3}\, r)) \\
        &= \ln(r^2) + 2 \ln(2r^2) + 3 \ln(3r^2) \\
        &= \ln(108 (r^2)^6) \\
        &= \ln\left(108 \left(\frac{x^2 + y^2 + z^2}{6}\right)^6\right) \\
        &= \ln\left(108 \left(\frac{a + b + c}{6}\right)^6\right)
    \end{align*}
    Da der Logarithmus streng monoton steigend ist, können wir aus $\ln(a b^2 c^3) \leq \ln\left(108 \left(\frac{a+b+c}{6}\right)^6\right)$ schlussfolgern:
    \[
        a b^2 c^3 \leq 108 \left(\frac{a + b + c}{6}\right)^6
    \]

\end{enumerate}

\end{document}
